\begin{afterword}
\summaryhead

The post asks how to admire something greatly without falling into a happy death spiral. It opens with Francis Bacon, whose Great Idea, the scientific method, promised enormous benefits and delivered them. So the remedy cannot be refusing to admire anything that much.

Nor, the author argues, can it be the standard criticisms of science, which probably came from people who wanted to undermine its authority; nor a selective search for negatives, which is rationalization; nor confining science to a narrow box, since science bears on nearly everything, even a parent's love.

The remedy is to treat each added nice claim as a burdensome detail and to ask for the specific chain of cause and effect behind it. A cancer patient will not be cured by publishing papers, whatever science may do for cancer. Claims that merely ``can't be disproven'' should add no happiness. Stuart Armstrong's advice, to split the Great Idea into independent parts and judge each separately, serves the same end. The post ends with a list of what does and does not work.

\responsehead

The post's method is to meet each new nice claim about a Great Idea with the question: by what specific chain of cause and effect would this be true? The problem is where the method is shown. The one realistic example of a false nice claim, that scientists need not take ethical responsibility because science will turn out well anyway, is dropped as ``not beyond dispute.'' In its place come claims that no admirer of science makes: publishing papers cures a cancer patient; a rule of $p < 0.05$ cures sociopaths. The post itself locates the danger in ``unsettled'' claims, like the example it dropped. On those the reader is not shown the method at work.

The remedies are offered without evidence. The closing list names five, and the post reports no test of any. The same was true of the remedy in ``Burdensome Details,'' which the post cites. Armstrong's advice has a further difficulty. The halo effect, as ``Affective Death Spirals'' defines it, acts between separate traits of one person, such as attractiveness and honesty. Splitting an idea into parts produces that setting. What would help is treating the parts ``as independent,'' which is the thing the halo makes hard. The post asserts that ``Three Great Ideas are far less likely to drive you mad than one Great Idea,'' and gives no evidence.

The examples are stated more strongly than the record. Bacon's claims of benefit came true, but his package also held a timetable: once nature's questions were answered, the invention of all causes and sciences ``would be the labor of but a few years.'' Four centuries later that work is not done. Judged in parts, as the post later advises, Bacon's Great Idea held up in some and not in others. The footnote says that Martine Rothblatt's search for a cure for her daughter's disease succeeded (``Successfully, I might add''). Her company made a treatment; pulmonary hypertension usually has no cure.

The post also attributes the standard criticisms of science, with ``Probably,'' to people whose pet beliefs science threatened. It names none, and some of the strongest early statements against the hydrogen bomb came from the scientists advising the US government in 1949.

In short: a method of asking for specifics, shown only on claims no admirer makes after the one realistic example is set aside, with remedies offered without evidence and examples stated more strongly than the record allows.
\end{afterword}
